% -----------------------------------------------------------------------------
% Scalar-only reduced momentum-space rules in strict SI units
% -----------------------------------------------------------------------------

LSZ 给出的物理场 Green 函数在标量散射中可整理为一组约化图因子。“约化”表示把每张图中重复出现的 SI 场归一化常数统一提出；四动量采用物理动量量纲，SI 归一化常数显式保留在 Green 函数与外线因子中。

质量为 $m$ 的实标量场，其动量-空间 Feynman 二点函数已由正则量子化得到
\begin{equation}
 \widetilde\Delta_F(p)
 =\frac{\ii\hbar^3c}{p^2-m^2c^2+\ii0^+}.
 \label{eq:scalar-physical-propagator-common-r15}
\end{equation}
分子 $\hbar^3c$ 来自采用的 SI 场归一化与 Fourier 约定；极点的位置 $p^2=m^2c^2$ 和 $+\ii0^+$ 边界条件与这个共同常数无关。因此定义标量约化内部线
\begin{equation}
 \boxed{
 \Delta_{\rm red}(p)
 \equiv\frac{\ii}{p^2-m^2c^2+\ii0^+}.}
 \label{eq:scalar-reduced-propagator-common-r15}
\end{equation}
它的 SI 量纲是物理动量的负二次方。这个定义不会把质量 $m$ 当成动量；分母中必须写 $m^2c^2$。


对实标量场 $\varphi$，把四次相互作用写成
\begin{equation}
 \boxed{
 \Lag_{\varphi^4}^{\rm int}
 =-\frac{g_4}{4!\,\hbar c}\varphi^4,}
 \qquad [g_4]=1.
 \label{eq:phi4-si-coupling-common-r15}
\end{equation}
由于 $[\varphi]=\sqrt{\mathrm{J/m}}$，有 $[\varphi^4]=\mathrm{J^2/m^2}$，而 $[\hbar c]=\mathrm{J\,m}$，所以右端确实具有拉格朗日密度的 SI 单位 $\mathrm{J/m^3}$。Dyson 插入与 LSZ 外腿截肢合并后，一个约化 $\varphi^4$ 顶角留下 $-\ii g_4$；这与 \eqnrefs{eq:phi4-lsz-result-r7} 的 SI 归一化完全一致。

